% \documentclass[12pt]{article}

%%v helpful MIT exam schema, see http://www-math.mit.edu/~psh/exam/examdoc.pdf
 \documentclass[8pt]{exam} 
% \qformat{\textbf{Question \thequestion}\quad \thepoints\hfill } % Large depth to make space}
 \renewcommand{\thesubpart}{(\roman{subpart})}
 \renewcommand{\subpartlabel}{\thesubpart}
 \footer{}{}{Page \thepage\ de \numpages}
\addpoints
% \pointsinmargin
% \boxedpoints
\renewcommand{\questionshook}{\setlength{\itemsep}{0.2cm}}
\renewcommand{\partshook}{
	\setlength{\itemsep}{0.1cm}
}

\usepackage{cmbright}

\pagestyle{empty} %for handouts
\addtolength{\jot}{2em} % this makes all formulae a bit more spaced vertically for handouts
\setlength{\parindent}{0pt} %we're not writing a novel
%\renewcommand{\arraystretch}{2} %makes tables taller, so easier to write in

\usepackage[fleqn]{amsmath}
\usepackage{amssymb} %standard classy maths symbols
\usepackage[a4paper, margin=1in]{geometry} %makes it easy to specify where to put stuff on the page

\usepackage{siunitx} %v handy way of getting SI units to typeset correctly...
\sisetup{locale = FR} %... and now it will use a comma for decimals

\usepackage[utf8]{inputenc}

\usepackage[french]{babel} % frenchify everything (e.g. a Table becomes a Tableau)
\usepackage{lmodern} % font with nicer French characters than standard
\usepackage{textcomp} % and more help for special characters

\usepackage{tcolorbox} % basic but good package for boxes around text
\tcbset{colback=black!5!white}
% \usepackage{color} % colours
%\newcommand{\fillin}{\color{black!1!white}} %makes the text disappear
%\definecolor{fillin}{rgb} {1.00,1.00,1.00}
%\renewcommand{\fillin}{\color{blue}} \definecolor{fillin}{rgb} {0.00,0.00,1.00} %makes the filling in text appear (in blue)

%\usepackage{enumitem} % can change the labels on lists, items, enumerates

\usepackage{tikz, pgfplots} % interval diagrams, sets, etc.
\usetikzlibrary{calc,trees,positioning,arrows,fit,shapes}

\begin{document}

\textbf{Nom:} \dotfill

\begin{center}
Let's study the function $g( x ) = x^2 -  2x - 3.$
\end{center}

\begin{questions}


\question[1] Give the y-intercept of $g$. \dotfill

\question[3] By Viète, find the zeros of $g$. \hfill $\Delta = b^2 - 4ac >0$ has 2 solutions $\frac{-b\pm\sqrt{\Delta}}{2a}$
\fillwithdottedlines{5cm}

\

\question[1] Complete the table of values.
\renewcommand{\arraystretch}{1.75}
\begin{tabular}{|l|l|l|l|l|l|l|l|}
\hline
$x$    & -2 & -1 & 0 & 1 & 2 & 3 & 4   \\ \hline
$g(x)$ &   \hspace{7mm}  &   \hspace{7mm}   & \hspace{7mm}   &   \hspace{7mm}   &  \hspace{7mm}  & \hspace{7mm}   &  \hspace{7mm}  \\ \hline
\end{tabular}

\vspace{4mm}

\begin{minipage}{0.5\textwidth}
\question[2] Plot the curve $y = g(x)$.

\begin{center}
\begin{tikzpicture}[scale=0.5]
\clip (-6.5,-6.2) rectangle (6.7,6.7);
\draw[help lines, step=1, white!50!gray] (-8.1,-6.6) grid (8.1,6.6);
%\draw[help lines,line width=.8pt,step=1] (-0.1,-0.1) grid (10.1,10.1);

\draw[thick, ->] (-6.1,0) -- (6.4,0) node[above] {$x$};
\draw[thick, ->] (0,-6.1) -- (0,6.4) node[right] {$y$};

\foreach \x in {-6, -5, ..., 6}  \draw[thick] (\x,0.1) -- (\x,-0.1) node[below] {$\x$};
\foreach \y in {-6, -5, ..., 6}  \draw[thick] (0.1,\y) -- (-0.1, \y) node[left] {$\y$};

%\draw[ultra thick] plot[smooth, domain=-8:8] (\x, 0.5*\x*\x-2*\x+1) node[above] {$f$};


\end{tikzpicture}
\end{center}
\end{minipage}
\begin{minipage}{0.5\textwidth}
\question[1] Find the x-intercepts (zeros) of $g$ again, but this time by \textbf{completing the square.}
\fillwithdottedlines{63mm}

\end{minipage}

\vfill

\textbf{6. Bonus.} Where would the graph of $y=g(x)$ cross the curve $y = 3x^2 + 6x + 5$ ? \hfill \emph{Work on the back.}

\end{questions}

\begin{tcolorbox}

\textbf{Nom de correcteur:}

\begin{center}
\gradetable[h][questions]
\end{center}

\end{tcolorbox}



\end{document}